\chapter{Analyse de sentiment et choix du modèle}
\label{chap:sentiment}

\section*{Introduction}
\addcontentsline{toc}{section}{Introduction}
L'analyse de sentiment est la partie d'intelligence applicative la plus visible dans Tastify. Elle ne remplace pas le jugement du gérant, mais elle transforme les avis textuels en indicateurs simples : sentiment positif, neutre ou négatif, score de confiance et modèle utilisé.

\section{Rôle de l'analyse de sentiment dans Tastify}
Dans Tastify, les avis clients donnent un retour direct sur l'expérience au restaurant. Un commentaire peut signaler un plat apprécié, un service lent, une attente trop longue, une mauvaise température ou une satisfaction générale. Sans traitement automatique, ces retours deviennent difficiles à suivre dès que leur nombre augmente.

L'analyse de sentiment joue donc deux rôles :

\begin{itemize}
  \item \textbf{opérationnel}, car elle aide le gérant à repérer rapidement les avis défavorables ;
  \item \textbf{analytique}, car elle alimente les statistiques de satisfaction affichées dans le tableau de bord.
\end{itemize}

L'implémentation couvre les statistiques de sentiment dans les indicateurs. Le résultat de l'analyse est donc exploité par l'application, il n'est pas seulement stocké.

\section{Données utilisées}
Le projet ne contient pas de jeu de données local destiné à entraîner un modèle propriétaire. Tastify utilise les avis saisis par les utilisateurs de l'application.

\begin{table}[H]
\centering
\caption{Données utilisées par le pipeline de sentiment}
\label{tab:donnees-sentiment}
\begin{tabularx}{\textwidth}{p{3.4cm}Y}
\toprule
\textbf{Champ} & \textbf{Utilisation} \\
\midrule
\pathcode{Avis.commentaire} & Texte libre analysé par le pipeline. C'est l'entrée principale du modèle. \\
\pathcode{Avis.note} & Note optionnelle entre 1 et 5. Elle est stockée, mais le code d'analyse actuel se base sur le commentaire. \\
\pathcode{Avis.lang_code} & Langue détectée simplement : \code{ar} si le texte contient des caractères arabes, \code{en} sinon. \\
\pathcode{AnalyseSentiment.label} & Résultat normalisé : \code{POSITIF}, \code{NEUTRE} ou \code{NEGATIF}. \\
\pathcode{AnalyseSentiment.score_brut} & Score de confiance renvoyé par Hugging Face ou par le fallback local. \\
\pathcode{AnalyseSentiment.modele_utilise} & Nom du modèle distant ou de la méthode locale réellement utilisée. \\
\bottomrule
\end{tabularx}
\end{table}

Le texte envoyé à Hugging Face est limité à 512 caractères dans la tâche Celery. Ce choix suffit pour des avis courts et évite d'envoyer des commentaires trop longs au modèle.

\section{Pipeline technique}
Le pipeline démarre lors de la création d'un avis dans le \code{AvisViewSet}. L'avis est enregistré rapidement, puis l'analyse est confiée à une tâche Celery.

\begin{figure}[H]
\centering
\fbox{\begin{minipage}{0.94\textwidth}
\centering
\schemaNode{Portail client\\saisie avis}
\schemaArrow
\schemaNode{API \code{/api/avis/}\\création}
\schemaArrow
\schemaNode{Tâche Celery\\analyse}

\vspace{0.45cm}

\schemaNode{Hugging Face\\modèle distant}
\schemaArrow
\schemaNode{Fallback local\\mots-clés}
\schemaArrow
\schemaNode{MySQL + Dashboard\\résultat et statistiques}
\end{minipage}}
\caption{Pipeline technique de l'analyse de sentiment}
\label{fig:pipeline-sentiment}
\end{figure}

La figure~\ref{fig:pipeline-sentiment} résume le flux : l'utilisateur soumet un avis, le backend l'enregistre, Celery lance l'analyse, puis le résultat est stocké et agrégé dans les statistiques.

\begin{figure}[H]
\centering
\includegraphics[width=0.95\textwidth]{capture_avis_sentiment.png}
\caption{Back-office : avis clients avec scores de sentiment}
\label{fig:capture-avis-sentiment}
\end{figure}

La figure~\ref{fig:capture-avis-sentiment} montre l'exploitation concrète du pipeline dans l'interface : chaque avis conserve son texte, sa note et un résultat de sentiment utilisable par le gérant pour prioriser les retours.

\section{Modèle réellement utilisé}
Le code du fichier \code{apps/avis/tasks.py} définit les modèles suivants :

\begin{itemize}
  \item \pathcode{nlptown/bert-base-multilingual-uncased-sentiment} pour les avis non arabes ;
  \item \pathcode{moussaKam/MARBERT-sentiment} pour les textes contenant des caractères arabes ;
  \item \pathcode{nlptown/bert-base-multilingual-uncased-sentiment} comme modèle de secours si le modèle arabe ne répond pas ;
  \item \pathcode{mots-cles-simple} comme secours local lorsque l'API Hugging Face n'est pas disponible ou lorsqu'aucun jeton n'est configuré.
\end{itemize}

Le modèle principal \pathcode{nlptown/bert-base-multilingual-uncased-sentiment} est un BERT multilingue affiné pour l'analyse de sentiment sur des avis produits en six langues, dont le français \cite{nlptown}. Il renvoie une prédiction sous forme d'étoiles de 1 à 5. Le code convertit ensuite cette sortie en trois classes : négatif, neutre ou positif. Ce choix correspond bien à un restaurant, car les commentaires clients ressemblent souvent à des avis courts sur un produit ou un service.

La base scientifique de ce type de modèle vient de BERT, qui apprend des représentations contextuelles bidirectionnelles et peut être adapté à des tâches de classification avec une couche de sortie \cite{bert}. Pour les textes arabes, le code utilise MARBERT, une famille de modèles conçus pour l'arabe et ses variantes dialectales \cite{marbert}.

\section{Conversion des sorties}
Les sorties du modèle ne sont pas exposées telles quelles au reste de l'application. Elles sont normalisées pour être lisibles dans le tableau de bord.

\begin{table}[H]
\centering
\caption{Normalisation des sorties de sentiment}
\label{tab:normalisation-sentiment}
\begin{tabularx}{\textwidth}{p{3.2cm}p{3.4cm}Y}
\toprule
\textbf{Sortie modèle} & \textbf{Label Tastify} & \textbf{Score métier} \\
\midrule
\code{5 stars}, \code{4 stars}, \code{POSITIVE}, \code{LABEL\_2} & \code{POSITIF} & $+15$, avis favorable. \\
\code{3 stars}, \code{NEUTRAL}, \code{LABEL\_1} & \code{NEUTRE} & $0$, signal sans polarité forte. \\
\code{1 star}, \code{2 stars}, \code{NEGATIVE}, \code{LABEL\_0} & \code{NEGATIF} & $-15$, avis défavorable à surveiller. \\
\bottomrule
\end{tabularx}
\end{table}

Cette conversion rend l'indicateur plus lisible pour le gérant, tout en gardant dans \code{score\_brut} la confiance renvoyée par le modèle.

\section{Comparaison avec des alternatives}
Le tableau~\ref{tab:comparaison-modeles} compare le modèle principal avec plusieurs alternatives possibles.

\begin{table}[p]
\centering
\scriptsize
\caption{Comparaison des approches de sentiment}
\label{tab:comparaison-modeles}
\begin{tabularx}{\textwidth}{p{3.2cm}p{2.5cm}YYY}
\toprule
\textbf{Approche} & \textbf{Statut dans Tastify} & \textbf{Forces} & \textbf{Limites} & \textbf{Pertinence} \\
\midrule
BERT multilingue NLP Town & Modèle principal & Supporte le français, l'anglais, l'espagnol, l'italien, l'allemand et le néerlandais ; entraîné sur des avis ; utilisable via API \cite{nlptown,hftextclassification}. & Moins spécialisé qu'un modèle français pur ; dépend d'un service externe si utilisé via API. & Très adapté aux avis courts de restaurant multilingues. \\
DistilCamemBERT sentiment & Alternative française & Spécialisé français ; modèle plus léger ; résultats documentés sur Amazon Reviews et Allociné \cite{distilcamembert,camembert}. & Ne couvre pas naturellement l'arabe ou l'anglais ; demande un routage par langue. & Bon choix si le périmètre devient surtout francophone. \\
XLM-R / XLM-T & Alternative multilingue & Bonne base de transfert multilingue ; XLM-R est conçu pour de nombreux langages \cite{xlmroberta}; XLM-T vise les textes sociaux multilingues \cite{xlmt}. & Demande un choix de modèle fine-tuné et des validations supplémentaires. & Intéressant pour une version multilingue plus avancée. \\
MARBERT sentiment & Modèle arabe utilisé & Appuyé sur la famille MARBERT, adaptée aux variétés arabes et au texte social \cite{marbert}. & Spécifique aux textes arabes ; ne remplace pas le modèle multilingue pour le français. & Adapté aux avis en arabe ou en dialecte marocain. \\
Mots-clés / TF-IDF local & Fallback partiel & Fonctionne hors ligne, transparent et peu coûteux ; le fallback actuel par mots-clés est déjà implémenté. & Faible face à l'ironie, aux négations et au vocabulaire varié ; TF-IDF n'est pas implémenté dans le dépôt. & Utile comme secours, mais insuffisant comme moteur principal. \\
\bottomrule
\end{tabularx}
\end{table}

\section{Justification du choix}
Le choix du modèle \pathcode{nlptown/bert-base-multilingual-uncased-sentiment} est défendable pour quatre raisons :

\begin{enumerate}
  \item \textbf{Proximité avec le domaine des avis.} Le modèle est affiné sur des avis, ce qui se rapproche des commentaires laissés sur des plats ou sur un restaurant.
  \item \textbf{Support du français.} La fiche du modèle indique un entraînement incluant le français, ce qui correspond au contexte du rapport et aux interfaces francophones.
  \item \textbf{Intégration rapide.} L'API Hugging Face permet de lancer l'inférence sans embarquer un modèle lourd dans le backend.
  \item \textbf{Continuité du service.} Le fallback local permet de produire un résultat même sans jeton Hugging Face ou en cas d'erreur réseau.
\end{enumerate}

Ce choix est cohérent avec le périmètre PFA et couvre les avis multilingues attendus dans Tastify. Pour une version de production, l'étape suivante serait de mesurer la qualité du modèle sur des avis du restaurant, constituer un petit jeu de validation local, puis comparer les scores entre le modèle multilingue, un modèle français spécialisé et un modèle arabe.

\section{Limites de la fonctionnalité}
Les limites suivantes sont identifiées :

\begin{itemize}
  \item la détection de langue reste simple : elle détecte l'arabe par présence de caractères arabes et classe le reste comme \code{en} ;
  \item le fallback local repose sur des listes de mots positifs et négatifs, pas sur un modèle statistique entraîné ;
  \item le projet ne fournit pas de métriques d'évaluation sur un corpus d'avis Tastify ;
  \item la note de 1 à 5 n'est pas combinée avec le commentaire pour améliorer la prédiction ;
  \item l'utilisation de Hugging Face suppose une configuration correcte du jeton \code{HUGGINGFACE\_API\_TOKEN}.
\end{itemize}

\section*{Conclusion}
\addcontentsline{toc}{section}{Conclusion}
L'analyse de sentiment de Tastify est intégrée au flux métier : elle tourne en arrière-plan, stocke son résultat et alimente les tableaux de bord. Le modèle principal correspond au contexte des avis clients, et une évaluation locale pourra renforcer encore sa précision dans une version de production.
